\chapterstart{表面分析技术基础}%
{本章导读：表面分析的共同前提是「信号必须在表面产生、并能量无损地逃出来」，所以先要回答的不是
「用哪台谱仪」，而是「信号来自多深、这个表面能保持多久」。本章讲表面的定义、表面与体相的差别、
信息深度与超高真空的关系（9.1），再以 X 射线光电子能谱（XPS）为样本讲清基本原理、仪器组成与
实验方式（9.2）。XPS 的定量、分峰与化学态判读，以及 AES、UPS、SIMS 等方法的选择，留待后续分片。}
\label{ch:09}

% ================================================================
%  第 9 章 表面分析技术基础（分片 a）
%  源：S2（下册）第 9 章，书内 p.567—614 / PDF p.59—106
%  覆盖：9.1 关于表面；9.2.1 历史进程（导航保留）；9.2.2 基本原理；
%        9.2.3 仪器与主要组成部件；9.2.4 实验方式
%  状态：分片 a 完稿；9.2.5 及以后由分片 b/c/d 续写
% ================================================================

\subsection{关于表面}

\begin{probout}
表面分析测量的是固体最外层约 0.5—10 nm 内的\textbf{元素组成、化学态、原子排列与电子结构}。
原始数据只有一条曲线：某类出射粒子（光电子、俄歇电子、二次离子、散射离子）的强度随能量（或
质量）的分布。由此换算出的量有三类：\textbf{元素种类与相对含量}（原子百分比）、\textbf{化学态}
（价态与化学环境），以及它们沿深度与横向的分布。它不直接给出体相组成，也不直接给出催化活性；
含量是相对值，其定量依赖相对灵敏度因子与「采样深度内组成均匀」这一假设。
\end{probout}

\subsubsection{表面的定义与厚度}

两种异态物质之间紧密接触的层称为界面，物体与真空或气体接触的界面称为表面，表面分析着重研究
固体表面，即气-固两相界面（特殊情况下也研究离子液体等特殊液态物质的表面）。表面层通常指一到
两个单层，信息来自零点几纳米到几纳米深处。这个尺度有多小：金属表面 $1$ cm$^{2}$ 的表面层区域
内约有 $10^{15}$ 个原子，而 $1$ cm$^{3}$ 立方体内的原子总数约 $10^{23}$ 个，表面原子只占体相
原子总数的 $10^{-8}$，即 $10^{-6}\%$——用体相灵敏的方法看表面，等于在约一亿个原子里找一个
原子。\tocheck{未对照原页面图像：原 OCR 作「百分数为 106\%」，上标负号丢失；按
$10^{15}/10^{23}=10^{-8}$ 重建。}

\subsubsection{表面与体相的差别}

表面是固体的终端，是晶体三维周期结构与真空之间的过渡区。在动力学与热力学平衡的前提下，表面的
化学组成、原子排列、原子振动状态均有别于体相，原因有两条：① 表面向外一侧没有邻近原子，存在
不饱和化学键即\textbf{悬挂键}，因此化学性质很活泼，极易吸附、氧化与重构；② 固体内部的三维
周期结构在表面中断，表面原子的\textbf{电子状态}（表面态、功函数、振动模式）也有异于体相。
所以表面分析的内容包括表面的形貌、化学组成、原子结构、原子态与电子态——每一项都需要不同的
探测粒子与不同的信息深度，没有一种技术能同时给出全部答案。
\src{S2 §9.1，书内 p.569—570 / PDF p.60—61}

\subsubsection{真实表面的保持：为什么必须超高真空}

在通常的大气环境下，一个新鲜表面很快就被大气组分沾污。\qty{25}{\celsius} 空气中 \ce{N2} 在
表面上形成单层所需的时间随压强急剧变化：$760$ torr 时只需约 $2.9\times10^{-9}$ s，
$10^{-6}$ torr（约 $1.3\times10^{-4}$ Pa）时约 $2.2$ s，$10^{-9}$ torr 时才延长到约
$2.2\times10^{3}$ s。换一种算法：在 $10^{-4}$ Pa、$300$ K 的 \ce{N2} 气氛中，每平方厘米每秒
有 $3.88\times10^{14}$ 个分子撞到表面，而典型固体表面约 $10^{15}$ 原子/cm$^{2}$，所以在
$10^{-4}$ Pa 这种「看起来已经很干净」的真空里，几秒钟就能盖满一个单层。表面被覆盖满一个单层
称为 1 Langmuir（1 L），$1$ L $\approx 10^{-4}$ Pa$\cdot$s，即 $10^{-4}$ Pa 下停留 1 s 就
等于 1 L 的暴露量。

真空区划对表面分析的意义并不相同：低真空（$10^{3}$—$10^{5}$ Pa）与中真空
（$10^{-1}$—$10^{3}$ Pa）下分子平均自由程 $\lambda$ 小于或接近容器直径 $d$，以分子之间的
碰撞为主，只够样品传输与粗处理；高真空（$10^{-6}$—$10^{-1}$ Pa）下 $\lambda>d$，但单层覆盖
时间仍只有秒级，只够快速进样；只有超高真空（$10^{-10}$—$10^{-6}$ Pa）把单层形成时间延长到
$10^{2}$—$10^{6}$ s，才是表面分析的常规工作区（极高真空 $<10^{-10}$ Pa，仅特殊需求使用）。
超高真空起两个作用：使一个较为真实的
表面在相当长的时间（$10^{2}$—$10^{6}$ s）内得以保持，以完成一次真正的表面分析；同时使出射
粒子（主要是电子）尽量少与气体分子碰撞，从而能量无损地到达能量分析器。所以说，没有超高真空
就没有表面科学。推论是：表面分析的每一炉实验都把时间当作消耗品，从清洁完成到录谱结束必须
控制在单层时间以内。\src{S2 §9.1，书内 p.570—571 / PDF p.61—62，表 9-1、表 9-2}

\tocheck{未对照原页面图像：本节 $10$ 的负指数在原 OCR 层普遍丢失上标负号，已按物理含义与
表 9-2 的区划重建，需与书内 p.570—571 原页核对。}

\subsubsection{信息深度：电子逃逸深度与非弹性散射平均自由程}

$1$ keV 量级（XPS 激发源的典型能量）的 X 射线能穿透约 1000 nm 而不显著衰减，而具有同样能量的
电子只能穿透约 10 nm。正是这个差别使 XPS 成为表面灵敏的技术：即使 X 射线在整个穿透深度内都
激发了光电子，也只有最外面约 10 nm 内产生、并且能能量无损地走到表面的电子才携带元素特征信息。
光电子在途中会遇到弹性散射（只改变方向）与非弹性散射（能量也受损失）；一旦发生非弹性碰撞，
电子失去原有的元素特征，成为谱图上抬高的连续本底。因此 XPS 采集到的出射光电子来自深度

\begin{equation}\label{eq:sampling-depth}
d=3\lambda\cos\theta
\end{equation}

式中 $\theta$ 为出射电子与样品面法线的夹角；$\lambda$ 为光电子的非弹性散射平均自由程
（\term{IMFP}{inelastic mean free path}{非弹性散射平均自由程}）。各深度的贡献按
$I(d)/I_{\infty}=1-\exp(-d/\lambda)$ 累积：$\lambda$ 处信号强度 63.2\%，$2\lambda$ 处
86.5\%，$3\lambda$ 处 95.0\%，$4\lambda$ 处 98.2\%，$5\lambda$ 处 99.3\%。一般把 XPS 的采样
深度定义为 $3\lambda$——出射光电子经过 $3\lambda$ 行程后只剩总量的 5\%。
\src{S2 §9.1—§9.2.2，书内 p.582—583 / PDF p.73—74，式(9-2)}

\begin{keycond}
\equref{eq:sampling-depth} 的成立条件：① $\lambda$ 沿出射路径恒定，即采样深度内组成与密度
均匀；② 只计非弹性散射造成的损失，忽略弹性散射对轨迹的随机化（它使实际信息深度略大于
$3\lambda\cos\theta$）；③ $\theta$ 是出射方向与\textbf{样品面法线}的夹角，$\theta=0$ 时采样
最深、$\theta$ 增大时采样变浅；④ 单晶表面的点阵衍射（通道效应）与表面粗糙度都会破坏这一关系。
量纲检查：$\lambda$ 与 $d$ 同为长度，$\cos\theta$ 无量纲，两侧量纲一致，式中不含经验常数，
可换用任何长度单位。另需注意「采样深度 $=3\lambda$」只是取「剩 5\% 信号」的工程约定，不是
物理边界：更深处产生的电子仍可能逃出，只是概率按指数衰减。
\end{keycond}

$\lambda$ 随电子动能变化，可用 Seah 与 Dench（1979）的经验式粗略算得：

\begin{equation}\label{eq:imfp}
\lambda=
\begin{cases}
538E^{-2}+0.41(\delta E)^{1/2}, & \text{纯元素}\\
2170E^{-2}+0.72(\delta E)^{1/2}, & \text{无机化合物}\\
49E^{-2}+0.11E^{1/2}, & \text{有机化合物}
\end{cases}
\end{equation}

式中 $E$ 为光电子能量（eV）；$\delta$ 为单原子层厚度（nm）；前两式的 $\lambda$ 以单层数为单位，
第三式以 mg/cm$^{2}$ 为单位。量纲检查：两式中的两项都必须是长度，所以 $538$、$2170$、$49$、
$0.41$ 等系数并非纯数、各自带单位（$538$ 的量纲是 nm$\cdot$eV$^{2}$）；把 $E$ 换成 keV 或把
$\delta$ 换成 Å 都必须重新标定系数。
\tocheck{未对照原页面图像：原 OCR 的无机化合物式作 $2170E^{-2}+0.72(\delta E)$，缺 $1/2$ 次方，
按 Seah—Dench 式的标准形式与量纲一致性重建为 $(\delta E)^{1/2}$；第三式 $\lambda$ 的单位原书
OCR 作「mg/cm$^{2}$」，而 Seah—Dench 原文常以 mg/m$^{2}$ 给出，单位口径存疑。}
\src{S2 §9.2.2，书内 p.583 / PDF p.74，式(9-3a—9-3c)}

\begin{figure}[htbp]
\centering
\begin{tikzpicture}[x=1mm,y=1mm,>=Stealth]
% 气相
\fill[gray!12] (0,0) rectangle (78,15);
% 固体
\fill[gray!25] (0,-26) rectangle (78,0);
% 采样深度带
\fill[gray!45,opacity=0.40] (0,-11) rectangle (78,0);
\draw[mainline] (0,0) -- (78,0);
\draw[dashed,thinline] (0,-11) -- (78,-11);
\draw[<->,thinline] (3.4,0) -- (3.4,-11)
  node[midway,right,inner sep=1.5pt] {$3\lambda$};
% 气相说明
\node[anchor=north west,text width=40mm,align=left,inner sep=1pt] at (2,14)
  {气相（$10^{-8}$ Pa 量级）：出射电子不与气体分子碰撞};
% 入射 X 射线
\draw[->,mainline] (62,14) -- (34,-8) node[midway,above right,inner sep=1pt] {$h\nu$};
% 法线与出射电子
\draw[dashed,thinline] (34,-8) -- (34,5);
\draw[->,mainline] (34,-8) -- (26,0)
  node[midway,above left,inner sep=1pt] {$E_{\mathrm{k}}$};
\draw[thinline] (34,-2) arc[start angle=90,end angle=135,radius=6];
\node[fill=white,inner sep=0.5pt] at (31.6,-2.2) {$\theta$};
% 深部电子：非弹性散射后无法逃出
\draw[->,thinline] (56,-24) -- (48,-16);
\draw[->,thinline,dashed] (48,-16) -- (41,-9);
\node[anchor=south west,text width=20mm,align=left,inner sep=1pt,fill=white]
  at (58,-25) {深部电子：非弹性散射失去能量，成为连续本底};
% 采样深度带说明
\node[anchor=south west,text width=34mm,align=left,inner sep=1pt] at (1,-25)
  {采样深度带：法向出射时 $3\lambda$ 内贡献约 95\% 信号};
\node[anchor=east,inner sep=1pt] at (76.5,-2.6) {固体};
\end{tikzpicture}
\caption{表面分析信号来源与逃逸深度示意。读者应看三处：① 入射 X 射线穿透约 1000 nm，远大于
电子的逃逸深度，信息深度由\textbf{电子}而非光子决定；② 只有采样深度带内产生、且沿出射方向未
发生非弹性散射的电子才形成弹性峰，其信息深度为 $d=3\lambda\cos\theta$；③ 更深处产生的电子经
非弹性散射后成为连续本底，本底随结合能升高而抬高。图中厚度与角度按可读性缩放，$\lambda$ 取
$3\lambda$ 带的 $1/3$。\schemamark\figsrc{据 S2 图 9-16，书内 p.582 / PDF p.73 改绘。}}
\label{fig:escape}
\end{figure}

$\lambda$ 与电子动能的关系不是单调的：约 50—100 eV 处取最小值（约 0.5 nm），向低能端与高能端
都上升。这带来一个实用推论：芯能级光电子的动能高（如 Al K$\alpha$ 激发的 C 1s 光电子约
1198 eV，$\lambda\approx1.1$ nm），采样深度较大；而俄歇电子的动能只有几百 eV，$\lambda$ 更小
（C KVV 约 270 eV 时 $\lambda\approx0.6$ nm），所以\textbf{俄歇峰的信息深度通常比同一元素的
芯能级光电子峰更浅}。

各种材料的采样深度粗略估计为：金属 $0.5$—$2$ nm，无机化合物 $1$—$3$ nm，有机物 $3$—$10$ nm。
还需记住：同一元素的不同能级、以及换用不同激发源，都会改变出射电子的动能，从而改变 $\lambda$
与采样深度——同一张谱上不同峰的表面灵敏度并不相同。
\src{S2 §9.2.2，书内 p.583 / PDF p.74，图 9-17}

\subsubsection{由信息深度决定的方法分工}

常见表面分析技术的原理与所获信息不同：\term{XPS}{X-ray photoelectron spectroscopy}{X 射线
光电子能谱}用 X 射线光子激发芯能级电子、测其动能以确定结合能，给出元素组成与化学态；
\term{AES}{Auger electron spectroscopy}{俄歇电子能谱}改用千电子伏电子束轰击，主要给元素组成
而束斑可以很小；\term{UPS}{ultra-violet photoelectron spectroscopy}{紫外光电子能谱}以紫外
光子激发价电子，给价带谱、功函数与费米能级信息；%
\term{SIMS}{secondary ion mass spectroscopy}{二次离子质谱}以高能离子束溅射、用质谱分析二次
离子，灵敏度极高且可测氢。XPS、UPS、AES 与 ISS 可在同一台谱仪中实现，共用一个分析器。
信息深度上 XPS、UPS、AES 为 $0.5$—$10$ nm，ISS 只有单层，SIMS 覆盖 $0.5$—$300$ nm，
AFM/STM 只给最上层原子的形貌、不给化学态；元素范围也不同（XPS/AES/UPS 为 Li—U，SIMS 为 H—U）。
应用时应记住两条原则：先弄清材料特性、所需信息与手段的能力和局限；不要奢望仅用一种技术满足
全部需求。\src{S2 §9.1，书内 p.569、p.571—573 / PDF p.60、p.62—64，表 9-4、图 9-1}

\begin{crosscheck}
表面敏感技术回答「最外层有什么、以什么价态存在」，不回答「体相是什么晶相、内部长什么样」：
后者要靠 \chapref{ch:04} 的 XRD 补体相晶相、靠 \chapref{ch:02} 的 TEM/EDS 补形貌与体相微区
成分、靠 \chapref{ch:05} 的化学吸附与程序升温方法补活性位与氧化还原行为。同为表面敏感的方法
之间也要互相回查：AES 束斑小、可扫描成像，但电子束损伤远大于 X 射线；UPS 给价带与功函数，
不给芯能级化学态；SIMS 灵敏度极高且可测氢，但定量困难、同样有溅射损伤；ISS 只给最外层原子。
定量结果还应用独立方法（XRF、ICP、EDS）的体相组成核对——两者的差异本身往往是信息（表面富集
或偏析），而不是哪一方「测错了」。
\end{crosscheck}

\subsection{X 射线光电子能谱（XPS）}

XPS 全称为 X-ray photoelectron spectroscopy，早期也称 ESCA（electron spectroscopy for
chemical analysis，化学分析用电子能谱），是用电子能谱仪测量 X 射线光子辐照样品表面所产生的
光电子与俄歇电子能量分布的方法。\src{S2 §9.2，书内 p.573 / PDF p.64}

\subsubsection{历史进程（导航保留）}

从 Hertz、Hallwachs 的观察（19 世纪 80 年代至 1888 年）到 Einstein 1905 年用能量量子化解释
光电效应，物理图像在 20 世纪初已经确立；1954 年 K. M. Siegbahn 制成世界上第一台光电子能谱
设备（Cu K$\alpha_{1}$ 源，$\lambda=0.15418$ nm，$h\nu\approx8042.6$ eV），精确测定了元素
周期表中各种原子的内层电子结合能。真正让它进入表面科学的是 20 世纪 60 年代的发现：硫代硫酸钠
中 \ce{S^{6+}} 与 \ce{S^{2-}} 的 S 2p 结合能相差约 6 eV，即\textbf{原子内层电子结合能的变化
可以提供分子结构与化学态的信息}——这是 XPS 的立命之本，K. M. Siegbahn 因此获 1981 年诺贝尔
物理学奖。仪器型号与人物沿革不再展开，需要时可回查原书图 9-2—图 9-8。
\src{S2 §9.2.1，书内 p.573—577 / PDF p.64—68}

\subsubsection{基本原理}

\paragraph{光电效应与基本方程}

\begin{principle}
XPS 的基本原理就是光电效应（光致发射、光电离）：样品被特征能量的光子辐照，测量出射光电子的
动能，从中获得所需信息。原子中不同能级上的电子具有不同的结合能；当一束能量为 $h\nu$ 的入射
光子与样品中的原子相互作用时，单个光子把全部能量交给原子中某壳层（能级）上一个受束缚的
电子。如果光子能量大于该电子的结合能，电子就脱离原来受束缚的能级，剩余能量转化为动能并以
一定动能从原子中发射出去成为自由电子，原子本身则成为激发态的离子。从激发到发射的过程极快
（约 $10^{-16}$ s），其能量守恒关系即爱因斯坦方程 \equref{eq:xps}。
\end{principle}

\begin{equation}\label{eq:xps}
E_{\mathrm{k}}=h\nu-E_{\mathrm{B}}-\phi_{\mathrm{sp}}
\end{equation}

式中 $E_{\mathrm{k}}$ 为出射光电子的动能，eV；$h\nu$ 为 X 射线源光子的能量，eV；
$E_{\mathrm{B}}$ 为特定原子轨道上的结合能，eV；$\phi_{\mathrm{sp}}$ 为谱仪的功函数，eV，主要由
谱仪材料与状态决定，对同一台谱仪基本是一个常数、与样品无关，平均值 $3$—$5$ eV。移项即得
$E_{\mathrm{B}}=h\nu-E_{\mathrm{k}}-\phi_{\mathrm{sp}}$：实测的是动能，报告的是结合能，中间的
$\phi_{\mathrm{sp}}$ 只能靠标定消去。\src{S2 §9.2.2，书内 p.579 / PDF p.70，式(9-1)}

\begin{keycond}
\equref{eq:xps} 的成立条件：① 样品与谱仪有良好电接触并共同接地，使二者费米能级处于同一能级
（导电样品，原书图 9-11）；绝缘样品做不到这一点，谱线会整体漂移（原书图 9-12）；② $h\nu$ 为
单一能量；③ $\phi_{\mathrm{sp}}$ 在同一台谱仪、同一工作状态下视为常数；④ 该式只描述单电子
电离的能量守恒，不含终态弛豫对峰位与峰形的修正。量纲检查：$h\nu$、$E_{\mathrm{B}}$、
$E_{\mathrm{k}}$、$\phi_{\mathrm{sp}}$ 全部是能量，取同一单位（eV）即可，式中没有需要标定的
系数；若改用 J 只需统一换算（$1$ eV $=1.602\times10^{-19}$ J）。极限情况：当
$h\nu<E_{\mathrm{B}}+\phi_{\mathrm{sp}}$ 时无光电子出射——这给出 XPS 可测能级的硬边界。
\end{keycond}

\paragraph{两条标定路线}由 \equref{eq:xps} 可见，$\phi_{\mathrm{sp}}$ 与样品荷电都会让整条谱线
一起平移，因此实测谱必须先标定，而且要把两件事分开。\textbf{谱仪能量标尺标定（仪器问题）}：
用 Au 4f$_{7/2}$、Cu 2p$_{3/2}$、Ag 3d$_{5/2}$ 等金属标样在同一台谱仪上测量，把已知结合能对准，
以消去 $\phi_{\mathrm{sp}}$ 与标尺的非线性。\textbf{样品荷电标定（样品问题）}：绝缘样品的荷电
使全部峰一起平移，必须用中和枪并把一个已知位置的峰调到指定值；最常用的是污染碳 C 1s 取
$284.8$ eV——它几乎无处不在、与样品处于同一荷电状态、一次录谱就能得到，代价是污染碳自身的
结合能随真空、样品与荷电状态在约 $284.5$—$285.0$ eV 之间漂移，用它定标等于把这个不确定度
转移给全部峰。两者不能互相替代：蒸镀 Au 做内标可辅助判断荷电，但只有 Au 与样品处于同一电位时
才成立。\tocheck{未对照原页面图像：Au 4f$_{7/2}$ 约 $84.0$ eV 与污染碳 C 1s $284.8$ eV 均为
通用标定值，分片 a 的源文本未给出这些数值，需与后续分片或原书核对后再作为书上给出的值引用。}

可被激发的能级还受 $h\nu$ 限制：常用 Mg K$\alpha$（$1253.6$ eV）与 Al K$\alpha$
（$1486.6$ eV）只能激发出结合能小于 $h\nu-\phi_{\mathrm{sp}}$ 的壳层。以金为例，Al K$\alpha$
下最深的可测能级是 4s，更内层的能级无法激发；H 与 He 没有内层电子，XPS 无法测量，所以 XPS 的
元素范围下限是 Li。

\paragraph{弛豫：X 射线荧光与俄歇电子}电离产生的终态离子不稳定、处于高激发态，会自发弛豫
（退激发）而变为稳态，弛豫分辐射弛豫与非辐射弛豫两种：前者发射 X 射线荧光，后者发射俄歇电子。
以 X 射线光子把能量转移给芯能级电子为例（原书图 9-15）：（a）从 $n$ 电子初态发生光电发射，
即光电子出射；（b）处于 $(n-1)$ 电子状态的原子通过把更高能级的电子退到空穴而重排；（c）重排
释放的能量可以交给另一个电子使其发射，这个电子就是俄歇电子，也可以以 X 射线光子的形式释放。
由此得到一条极有用的判别依据：\textbf{光电子峰的表观位置随 $h\nu$ 移动，而俄歇峰的动能由原子
自身能级决定、与 $h\nu$ 无关}。换靶复测（例如 Mg K$\alpha$ 换 Al K$\alpha$）时，位置不动的是
光电子峰，位置移动的是俄歇峰——这既是识别未知峰的判据，也是排除 X 射线卫星线的判据
（图~\ref{fig:xpsspec}）。\src{S2 §9.2.2，书内 p.581—582 / PDF p.72—73，图 9-15}

\begin{figure}[tbp]
\centering
\begin{tikzpicture}
\begin{groupplot}[
  group style={group size=2 by 1, horizontal sep=8mm},
  width=78mm, height=46mm,
  x dir=reverse, ymin=0, ytick=\empty, scaled ticks=false,
]
\nextgroupplot[xmin=0, xmax=1000, ymax=2.6, xtick={0,200,400,600,800,1000},
  xlabel={结合能 $E_{\mathrm{B}}$ / eV}, ylabel={强度（任意单位）}]
\addplot[domain=0:1000,samples=300,mainline,black]{0.30*exp(-x/700)
+1.00/(1+((x-285)/6)^2)+0.35/(1+((x-287.5)/6)^2)+0.12/(1+((x-291.5)/6)^2)
+0.55/(1+((x-531)/7)^2)+0.30/(1+((x-983.6)/8)^2)
+0.21*(1+(x-285)/sqrt((x-285)^2+9))+0.05*(1+(x-287.5)/sqrt((x-287.5)^2+9))
+0.02*(1+(x-291.5)/sqrt((x-291.5)^2+9))+0.12*(1+(x-531)/sqrt((x-531)^2+9))
+0.08*(1+(x-983.6)/sqrt((x-983.6)^2+9))};
\addplot[domain=0:1000,samples=300,thinline,black,dashed]{0.12+0.30*exp(-x/700)
+1.00/(1+((x-285)/6)^2)+0.35/(1+((x-287.5)/6)^2)+0.12/(1+((x-291.5)/6)^2)
+0.55/(1+((x-531)/7)^2)+0.30/(1+((x-750.6)/8)^2)
+0.21*(1+(x-285)/sqrt((x-285)^2+9))+0.05*(1+(x-287.5)/sqrt((x-287.5)^2+9))
+0.02*(1+(x-291.5)/sqrt((x-291.5)^2+9))+0.12*(1+(x-531)/sqrt((x-531)^2+9))
+0.08*(1+(x-750.6)/sqrt((x-750.6)^2+9))};
\draw[dashed,thinline] (axis cs:983.6,0) -- (axis cs:983.6,1.80);
\draw[dashed,thinline] (axis cs:750.6,0) -- (axis cs:750.6,1.80);
\node[anchor=north west,text width=74mm,align=left,inner sep=1pt] at (axis cs:1000,2.55)
  {实线：Al K$\alpha$（1486.6 eV）；虚线：Mg K$\alpha$（1253.6 eV，为便于分辨整体上移 0.12）。
  换靶后 285.0（C 1s）与 531（O 1s）两个弹性峰位置不动，俄歇 O KVV 峰则由 983.6 移到
  750.6 eV。};
\nextgroupplot[xmin=274, xmax=296, ymax=1.35, xtick={276,280,284,288,292},
  xlabel={结合能 $E_{\mathrm{B}}$ / eV}]
\addplot[domain=274:296,samples=250,mainline,black]{0.18+0.008*(x-274)
+0.85/(1+((x-285)/0.55)^2)+0.45/(1+((x-286.8)/0.6)^2)+0.25/(1+((x-288.5)/0.65)^2)
+0.15/(1+((x-291.5)/0.9)^2)
+0.10*(1+(x-285)/sqrt((x-285)^2+0.09))+0.05*(1+(x-286.8)/sqrt((x-286.8)^2+0.09))
+0.03*(1+(x-288.5)/sqrt((x-288.5)^2+0.09))};
\draw[dashed,thinline] (axis cs:285,0.02) -- (axis cs:285,0.72);
\draw[dashed,thinline] (axis cs:286.8,0.02) -- (axis cs:286.8,0.72);
\draw[dashed,thinline] (axis cs:288.5,0.02) -- (axis cs:288.5,0.72);
\draw[dashed,thinline] (axis cs:291.5,0.02) -- (axis cs:291.5,0.72);
\node[anchor=north west,text width=33mm,align=left,inner sep=1pt] at (axis cs:283.4,1.33)
  {C 1s 分峰（示意）：$285.0$ C--C/C--H；$286.8$ C--O；$288.5$ C=O；
  $291.5$ 伴峰（shake-up，$\pi\rightarrow\pi^{*}$）};
\end{groupplot}
\end{tikzpicture}
\caption{模型计算的 XPS 谱示意。(a) 全谱：同一含碳含氧样品在 Al K$\alpha$（实线）与
Mg K$\alpha$（虚线，为便于分辨整体上移 0.12）激发下的谱形，读者应看三处：① 连续本底随结合能
升高而抬高，来自非弹性散射电子（成因见图~\ref{fig:escape}）；② 弹性峰（C 1s 285.0 eV、
O 1s 531 eV）在换靶后峰位不动；③ 俄歇 O KVV 峰的表观结合能由 983.6 eV 变为 750.6 eV，
即随激发光子能量移动。(b) C 1s 窄区放大：主峰 285.0 eV，化学位移峰 286.8 与 288.5 eV，
高结合能侧 291.5 eV 的 $\pi\rightarrow\pi^{*}$ 伴峰（shake-up）。(a) 的横轴按 XPS 惯例由右向左
增大。峰形取洛伦兹型 $A/[1+((E_{\mathrm{B}}-p)/w)^2]$（$w=6$—8 eV，为可读性放大），本底为
指数衰减项 $B_{0}e^{-E_{\mathrm{B}}/700}$ 加各峰下的阶跃项；峰位、峰宽与强度均为示意设定，
不代表任何实测样品。\modelmark}
\label{fig:xpsspec}
\end{figure}

\paragraph{能级标记与自旋轨道分裂}原子中单个电子的运动状态用量子数 $n,l,m_{l},m_{s}$ 描述：
$n$ 为主量子数，能量主要取决于它，$n=1,2,3,4,\dots$ 对应 K、L、M、N 壳层；$l$ 为角量子数，
$l=0,1,2,3,\dots$ 对应 s、p、d、f 能级；$m_{l}$ 决定电子云取向；$m_{s}$ 为自旋量子数，取
$\pm1/2$。轨道运动与自旋运动之间的耦合使能级分裂，当 $l>0$ 时用内量子数 $j$ 表示：

\begin{equation}\label{eq:jsplit}
j=|l\pm1/2|
\end{equation}

所以 $l=0$ 时 $j=1/2$（s 能级不分裂），$l>0$ 时 $j=l+1/2$ 或 $l-1/2$，在谱图上表现为双峰，
如 Ag 3d$_{5/2}$ 与 3d$_{3/2}$；s 能级的内量子数 $1/2$ 通常省略，如 C 1s 只有一个峰。XPS 一般
采集芯（内）能级谱（C 1s、Al 2p、Ag 3d、Au 4f 等），因为这些谱线强度强、分辨好。双峰的面积比
由简并度 $2j+1$ 决定，即 $p$ 为 $2\!:\!1$、$d$ 为 $3\!:\!2$、$f$ 为 $4\!:\!3$——这是判断
双峰是否成对、以及是否有物种被漏掉的独立约束。
\src{S2 §9.2.2，书内 p.584—585 / PDF p.75—76，表 9-5、图 9-18}

\paragraph{光致电离截面与灵敏度因子}光电效应的概率用光致电离截面 $\sigma$ 表示，定义为某能级
的电子对入射光子的有效能量转移面积，或一定能量的光子从某个能级激发出一个光电子的概率。入射
光子能量一定时有如下规律：① 同一原子中半径越小的壳层 $\sigma$ 越大；② 电子结合能与入射光子
能量越接近 $\sigma$ 越大；③ 不同原子中同一壳层，原子序数越大 $\sigma$ 越大。$\sigma$ 越大，
该能级的光电子峰强度就越大。常用的数据库有两个：Scofield 理论光致电离截面（基于 C 1s $=1$）
与 Wagner 实验灵敏度因子（基于 F 1s $=1$），各谱仪厂商的因子库都源自这两个数据源。这决定了
一条纪律：\textbf{定量时必须声明用了哪一套因子、基准是什么}，两套混用会带来系统性偏差。
XPS 的表面检测灵敏度可达 $10^{-3}$ 原子单层，但对体相检测的灵敏度只有 0.1\% 左右。
\src{S2 §9.2.2，书内 p.585—586 / PDF p.76—77，表 9-6、图 9-19}
\tocheck{未对照原页面图像：原 OCR 作「可以达到 103 原子单层」，上标负号丢失；按通行表述重建为
$10^{-3}$ 原子单层，需与书内 p.586 原页核对。}

\paragraph{从谱图到结论}

\begin{dataflow}
一条完整的判读链是：录谱（全谱定元素，目标元素窄谱定化学态）$\rightarrow$ 校准（谱仪标尺用
Au 4f$_{7/2}$ 等标样，样品荷电用内标或中和枪）$\rightarrow$ 扣背景（Shirley 或线性，扣法直接
改变峰面积）$\rightarrow$ 分峰拟合（约束峰形、半高宽上限、双峰间距与面积比）$\rightarrow$ 用
相对灵敏度因子把峰面积换算为原子百分比 $\rightarrow$ 与参考态（同批标准样品或数据库）比对给出
化学态。每一步都留下假设：标定物必须与样品处于同一荷电状态；背景扣法影响面积从而影响定量；
拟合约束必须写明，否则分出的物种不可复核；定量只给相对含量，且以采样深度内组成均匀为前提。
\end{dataflow}

\begin{obs}
\pedset{教学设定数值}某 Ni 基样品的 Ni 2p$_{3/2}$ 主峰实测结合能为 $855.4$ eV，在高结合能侧
约 $+6$ eV 处有一个宽的卫星峰；同一谱中污染碳 C 1s 位于 $284.8$ eV。样品在空气中放置一周后
重复测量，主峰移到 $856.2$ eV，而卫星峰与主峰的间距不变、峰形也无明显变化。
\end{obs}

\begin{infer}
峰位移动而双峰间距、卫星峰结构与峰形不变，说明这次移动来自荷电或标定条件的整体变化，而不是
Ni 的价态改变。若要把 $855.4$ eV 归属为某一价态，前提是：① 谱仪标尺已用金属标样校准；
② 荷电校准所用内标与样品同电位；③ 与同批参考态（金属 Ni、NiO 等）在同一条件下测得并对比。
三条缺一，峰位就只能写成「与某参考态接近」，不能作为价态的唯一证明。
\end{infer}

\subsubsection{仪器与主要组成部件}

XPS 仪器由配套的高、超高真空系统，电子能量分析器与探测器，激发源（XPS 用 Mg/Al 双阳极或
单色化 Al K$\alpha$ 源，AES 用电子枪，UPS 用紫外光源，ISS 用离子源），低能离子源，电子中和枪，
多自由度样品架，以及仪器控制与数据采集处理系统等组成。这些部件不是可选项的堆叠，而是围绕
「让少量低能电子在无场、无碰撞、无荷电干扰的条件下被计数」这一目标互相约束的整体
（表~\ref{tab:xpsparts}）。\src{S2 §9.2.3，书内 p.586—605 / PDF p.77—96，图 9-20、图 9-21}

\begin{table}[tbp]
\centering
\caption{XPS 谱仪的主要组成部件、常见形式与使用要点}
\label{tab:xpsparts}
{\srcsize
\begin{tabularx}{\twocolwidth}{@{}l p{56mm} X@{}}
\toprule
部件 & 常见形式与典型取值 & 使用要点 \\
\midrule
真空系统 & 快速进样室（$10^{-3}$ Pa 量级）、样品制备室、分析室（$10^{-8}$ Pa 量级）；旋片式
机械泵、涡轮分子泵组、溅射离子泵、钛升华泵 & 三室之间用大口径隔离阀与传输机构；分析室腔体用
$\mu$ 合金（Fe18/Ni75/Cr2/Cu5）屏蔽外磁场；压强低于设定值才允许开启各激发源 \\
X 射线源 & Mg/Al 双阳极：Mg K$\alpha$ $1253.6$ eV、线宽 $0.7$ eV，Al K$\alpha$
$1486.6$ eV、线宽 $0.85$ eV；单色化 Al K$\alpha$ 源（$\alpha$-石英 $10\bar{1}0$ 面，罗兰圆
半径常取 $250$ mm） & 非单色源带 K$\alpha_{3}$、K$\alpha_{4}$、K$\beta$ 卫星线与强的韧致辐射
本底，常用约 $2$ $\mu$m 的 Al 箔滤去大于 $1.5$ keV 的连续背景；单色化后线宽可降至约
$0.25$ eV，并可得小束斑，便于小面积、线扫描与成像 \\
电子能量分析器 & 静电场型半球 $180^{\circ}$ 分析器（SCA），常见直径 $220$—$450$ mm；外壳
双层 $\mu$ 合金 & 同心半球加 $\Delta V$（外半球为负），只有一定能量范围的电子能到达出口；
XPS 用 FAT/CAE 模式（绝对分辨率恒定、检测效率一致，适合面积法定量），AES 与 UPS 可用
FRR/CRR 模式 \\
透镜与传输 & 收集（预减速）透镜、多级静电透镜，可切换减速比 & 分辨能力与传输特性互相矛盾。
预减速后动能由 $E_{0}$ 降到 $E_{0}'$，分析器只需满足 $R_{b}'=(E_{0}/E_{0}')R_{b}$，系统
分辨率不变而分析器尺寸不必加大；同时拉开样品与分析器距离，减少次级电子本底 \\
探测器 & 单通道电子倍增器（增益 $10^{6}$—$10^{9}$）、$3/6/7/9$ 管阵列、微通道板（MCP），
多至 256 通道 & 待测电子流只有 $10^{-11}$—$10^{-8}$ A；阵列沿能量色散方向排列，输出按能量
移位后相加可提高灵敏度；MCP 单通道增益远低于单通道倍增器，不能只按通道数比较灵敏度 \\
荷电中和 & 早期为单束低能电子枪，现在多用低能 \ce{Ar+} 与低能电子束的双束中和 & 作用是减小
绝缘样品表面荷电，使作为内标的峰回到指定位置；还须消除荷电不均匀，否则会出现多重裂分谱峰，
被误读成「多种化学态」 \\
离子枪与样品架 & 差分抽气细聚焦扫描型（置于分析室）、大束斑高速型（置于制备室）；\ce{Ar}
团簇离子枪（团簇含 $2000$—$3600$ 个 \ce{Ar} 原子）；多自由度样品架 & 单 \ce{Ar+} 束轰击会
破坏表面状态，对有机物与聚合物尤甚；团簇源把初始动能分摊到上千个原子上，冲击能量低、损伤小。
变角分析要求样品面与样品架 $x$ 轴轴心共面 \\
\bottomrule
\end{tabularx}}
\figsrc{据 S2 §9.2.3，书内 p.586—605 / PDF p.77—96，图 9-20—图 9-42 整理。}
\end{table}

\paragraph{四点必要说明}① \textbf{真空}：三室分别负责快速进样、样品预处理与主分析，把热处理、
蒸镀、刻蚀等「脏活」放在制备室是为了保住分析室的 $10^{-8}$ Pa；泵有层级（旋片泵作前级，
涡轮分子泵或溅射离子泵作主泵，钛升华泵辅助），且只有压强低于设定值时联锁才允许开启激发源，
以保护灯丝、阳极靶与样品。② \textbf{分析器模式}：FAT/CAE 把进入分析器的电子减速到固定通能，
因 $\Delta E/E$ 中的 $E$ 固定而整条标尺上绝对分辨率不变，通能越低分辨越好但信号越弱（窄谱
$5$—$25$ eV、全谱 $100$—$200$ eV）；FRR/CRR 的 $\Delta E/E$ 恒定，但通能随能量变化、检测
效率不同，不利于定量，故 XPS 不用，而适合 AES 的峰对峰高度法与 UPS。③ \textbf{X 射线源}：
高能电子束撞击阳极产生内空穴并发射荧光 X 射线与韧致辐射，特征线能量固定；双阳极源通量高，
但带卫星线、次级电子本底、样品降解与灯丝沾污，故常用 Al 箔滤去大于 $1.5$ keV 的连续背景；
更好的办法是单色化——$\alpha$-石英 $10\bar{1}0$ 面间距 $0.425$ nm，对波长 $0.834$ nm 的
Al K$\alpha$ 线在 $78.5^{\circ}$ 处满足布拉格方程

\begin{equation}\label{eq:bragg-ch09}
n\lambda=2d\sin\theta
\end{equation}

数值核对：$2\times0.425\times\sin78.5^{\circ}=0.833$ nm，与 $0.834$ nm 相差约 0.1\%，自洽。
单色化可把 Al K$\alpha$ 线宽从 $0.85$ eV 降至约 $0.25$ eV，消除 K$\alpha_{3}$、K$\alpha_{4}$
伴峰与韧致辐射本底，并取得小束斑。④ \textbf{探测与中和}：待测电子流只有 $10^{-11}$—$10^{-8}$ A，
靠电子倍增器计数；中和枪减小绝缘样品荷电、使内标峰回到指定位置，并消除荷电不均匀（否则多重
裂分谱峰会被误读为多种化学态）。
\src{S2 §9.2.3，书内 p.588—604 / PDF p.79—95，图 9-21—图 9-42}
\tocheck{未对照原页面图像：原书（书内 p.604）称绝缘样品表面「负电荷富集」，而按光电发射过程，
绝缘样品通常因电子出射而带正电，这也是中和枪要提供低能电子的原因；两种说法方向相反。分片 a
只保留「中和不足或过度都会留下荷电」这一结论，不判断原书是否有误，请以原页面图像复核。}

\subsubsection{XPS 实验方式}

传统 XPS 给出固体表面某一小区域的元素组成、化学态与元素相对含量；现代 XPS 还能提供元素及其
化学态的横向与纵向（深度）分布，即线扫描、深度剖析与成像。常规点分析可使用 Mg-Al 双阳极或
单色化 X 射线源；用单色化源再配合分析器透镜部分的光阑，可实现小面积分析，目前最小采样面积
直径可小至 $9$ $\mu$m。

\paragraph{深度剖面：破坏性与非破坏性两条路线}破坏性方法用于获取深于 10 nm 的组分变化信息，
非破坏性方法只能给出表面 $1$—$10$ nm 内的组成变化。

\emph{离子刻蚀（破坏性）}先用常规方法分析原始表面，再以几百 eV 至几千 eV 的离子束刻蚀一定
时间、停止刻蚀并测谱，如此循环至所需深度。聚焦离子束斜向作用于表面，刻蚀坑不可能均匀，
特别是坑的边缘会使深度数据产生偏差，故要求刻蚀区域大于 XPS 采集出射光电子的范围并使用可旋转
样品托；最理想的刻蚀面积是测试面积的 100 倍，以确保测试区处于坑中心的平坦位置。刻蚀后的新鲜
表面极其敏感，须用高纯或超纯 \ce{Ar} 并注意系统残余气氛，否则新鲜表面的化学组成会被改变。
刻蚀深度用已知准确厚度的薄膜标样标定：常用的 \ce{Ta2O5}/Ta 标样由阳极化工艺制得、层厚控制
精确，例如 $30$ nm 标样的氧化层厚度就是 O 1s 与 Ta 4f$_{7/2}$ 强度-溅射时间曲线中Ⅰ段、Ⅱ段
与Ⅲ段 50\% 之和在横坐标上的投影段。副效应必须逐条记住：\textbf{择优溅射}（各元素溅射系数
不同，合金的新表面组成偏离原始组成）、热效应、离子注入、\textbf{表面局部被还原}（硝酸盐、
磷酸盐、碳酸盐在 $1$—$3$ keV \ce{Ar+} 轰击下可转化为氧化物，高价态氧化物最易被还原）、粗糙度
增加、表层重构（聚合物尤明显）与污染分析室。因此刻蚀应选低能离子束，大面积、长时间的清洁尽量
在制备室进行；对有机物与聚合物宜改用 \ce{Ar} 团簇离子源（团簇含 $2000$—$3600$ 个 \ce{Ar}
原子），初始动能分摊到上千个离子上，冲击能量低、化学破坏小。

\emph{变角 XPS（非破坏性）}XPS 测得的表面信息其实是表面与近表面多层信息的叠加，来自各层的
信息不同，这正是变角深度分析的基本概念：改变样品面的倾斜角度即精确改变出射角，利用采样深度与
出射角的关系得到元素浓度与深度的关系，

\begin{equation}\label{eq:arxps}
d^{*}=d\cos\theta=3\lambda\cos\theta
\end{equation}

式中 $\theta$ 为出射角（进入分析器方向的出射电子与样品法线的夹角），$\lambda$ 为非弹性散射
平均自由程；\equref{eq:arxps} 与 \equref{eq:sampling-depth} 是同一关系，只是写成随角度变化的
采样深度 $d^{*}$。$\theta=0^{\circ}$ 时采样最深，增大 $\theta$ 可获得更多表面层信息，
$\theta=85^{\circ}$ 时表面灵敏度提高约 10 倍（$\cos85^{\circ}\approx0.087$）。变角 XPS
（\term{ARXPS}{angle-resolved XPS}{变角 XPS}）只适用于表面层很薄（$1$—$10$ nm）的体系，
使用时还须注意：单晶表面的点阵衍射效应、表面粗糙度、样品面与样品架 $x$ 轴轴心共面、绝缘样品
荷电补偿随角度的变化（同样适用于 \ce{Ar+} 刻蚀剖面）、以及录谱范围要能容纳化学位移幅度。
与此相关的一种新方法是平行探测（结合位敏探测器）：由传输透镜、半球分析器与二维探测器
（微通道板）结合，在一个平面上收集能量色散信息、在另一平面收集角分布信息，不转动样品即可获得
大于 $60^{\circ}$ 范围的角度分布数据，得到无损的浅层化学态深度分布，且分析面积与荷电补偿不变。

\paragraph{线扫描与成像}用小束斑单色化 X 射线源，让细聚焦 X 射线束在样品表面线性移动（或束
固定、样品移动），可测得元素及其化学态的线强度分布；让束作 X-Y 扫描即得分布图。成像方式主要有
四种：扫描阳极使细 X 光束在样品上扫描、计算机控制样品台作 X-Y 扫描、在分析器透镜中间加静电
偏转板作光电子束扫描，以及平行（照相式）成像——后者每个像素都含一个完整谱，峰拟合后可再现
化学态的空间分布。
\src{S2 §9.2.4，书内 p.605—614 / PDF p.96—105，图 9-43—图 9-56}

\begin{measure}
录谱与校准的可操作条件：① \textbf{通能}：高分辨窄谱 $5$—$25$ eV，全谱 $100$—$200$ eV；
模式选 FAT/CAE，因为它的绝对分辨率在整条标尺上恒定、检测效率一致，可作面积法定量。
② \textbf{校准分两步}：先用 Au 4f$_{7/2}$、Cu 2p$_{3/2}$、Ag 3d$_{5/2}$ 等金属标样校准谱仪
能量标尺，再处理样品荷电——绝缘样品用中和枪并把内标（污染碳 C 1s 取 $284.8$ eV）调到指定
位置，且标定物与样品必须处于同一荷电状态。③ \textbf{深度剖析}：离子能量尽量低；刻蚀面积
不小于测试面积的 100 倍并用可旋转样品托；用高纯 \ce{Ar}；刻蚀速率用 \ce{Ta2O5}/Ta 标准样
标定；有机物与聚合物改用团簇离子源。④ \textbf{变角分析}：表面层厚度小于 10 nm；样品面与
样品架 $x$ 轴轴心共面；绝缘样品要检查荷电补偿随角度的变化。⑤ \textbf{定量}：同一批数据只用
一套相对灵敏度因子（Scofield 基准 C 1s $=1$，或 Wagner 基准 F 1s $=1$）并写明基准。
\end{measure}

\begin{pitfall}
\textbf{① 荷电效应冒充化学位移。}绝缘样品上，X 射线激发的电子发射与中和枪注入的电子并不总能
平衡，表面会留下净电荷。由 \equref{eq:xps} 可直接看出：表面带正电时出射电子被减速、实测动能
偏低，表观结合能偏高；中和过度则表观结合能偏低。判别办法不是看单个峰而是看整体：荷电使
\textbf{全部元素的峰一起平移}，且平移量随 X 射线通量、采集时间或中和条件改变；真实的化学位移
只影响价态发生变化的那些元素，且自旋轨道分裂间距、卫星峰结构与峰形会随峰位一起变化，并在
改变通量或中和条件后可以重现。荷电不均匀还会造成峰形展宽甚至多重裂分，把「一种化学态」读成
「多种化学态」——这时应先怀疑中和，而不是先加分峰。

\textbf{② X 射线还原样品。}非单色源的高能韧致辐射、次级电子与热量都可能使样品降解：高价态
氧化物、硝酸盐、磷酸盐、碳酸盐最易被还原（\ce{Ar+} 刻蚀有同样的还原作用）。判据是\textbf{时间
序列}：若连续录谱中高价态峰单调减弱、低价态峰单调增强，首先应怀疑辐照损伤，而不是「发现了新
物种」；降低通量、缩短单次采集、改用单色化源后变化减小，即可确认。

\textbf{③ 污染碳干扰。}真空残余的碳氢化合物在 X 射线或电子束下裂解沉积到表面，表现为 C 1s
随时间增长（常伴 O 1s）。危害有两层：若样品本身含碳，污染碳会掩盖真实碳物种；若用污染碳作
内标，则把它自身约 $284.5$—$285.0$ eV 的不确定度转移给全部峰。判别办法是比较新解理或刚刻蚀
完的表面与放置一段时间后的 C 1s 强度。

\textbf{④ 深度不均匀导致定量偏差。}相对灵敏度因子定量以「采样深度内组成均匀」为前提。核壳
结构、表面富集或偏析、刻蚀坑边缘的信号混合都会破坏这个前提：前者使表面组成不能代表整体，
后者把界面「展宽」成渐变层。自检办法有两个：做变角测量，峰强度比随出射角变化即说明深度不均
匀；或用同一元素动能差别较大的两个峰分别定量，结果不一致同样说明不均匀。此时定量只能表述为
「在 $3\lambda$ 深度内的相对含量」。
\end{pitfall}
